\subsubsection{Complete reconstruction of the CUE--Mellin peeling programme}
\label{sec:cue-mellin-complete-programme}

This subsection reconstructs the complete mathematical content of the three modern local
CUE--Mellin TeX versions.  Their exact source identities are:
\begin{quote}\small
maintained source (1,591 lines):\\
\texttt{4d4ed40c9968cc8cdfd28ebe129316ee04abcef47cef6683c01c47bd972daa4f}\\
predecessor (1,032 lines):\\
\texttt{21e450e24691c845c2d7d60b74e43555a854ccaf33407cbbe56d7dcc5a46b4a9}\\
predecessor (877 lines):\\
\texttt{7666d8a685914f91f0de96cb5e420c989001150a30a910bc33bb82249d61acea}\\
separate 17-page PDF:\\
\texttt{78eb0f5a545fa5c05b459aed0f973ae6fedcccd5615f84a914a5cd8179f47441}.
\end{quote}
They are versions of one programme, but deletion from a later version is not treated as a
mathematical retraction.  In particular, the earlier auxiliary Gauss law, ratio-lattice
count, and cubic trace calculation are audited below rather than lost.  The PDF was compared
through native PDF text with a fresh two-pass compilation of the maintained TeX.  The
mathematical body agrees; the 28 one-sided comparison lines are the two versions of fourteen
table-of-contents, pagination, and date lines.  Thus it is a rendered snapshot, not an
additional mathematical source.  No OCR is involved.

The human-source boundary is fixed first.  Haar CUE, Weyl integration, and the derivative
moment problem come from the published random-matrix literature
\cite{forrester2022,winn2012,conreyRubinsteinSnaith2006,dehaye2008,
groverMezzadriSimm2026,simmWei2026}.  The exact independent Verblunsky variables come from
Killip--Nenciu, with the derivative telescope derived in
\cref{thm:verblunsky-peel-derived} from their coordinates
\cite{killipNenciu2004}.  Bernstein's original construction supplies the approximation
operator used in the measure recovery \cite{bernstein1912}; Hausdorff is the historical
source for the finite-difference moment criterion \cite{hausdorff1921}.  The local formulas
below are then either consequences proved here from those inputs or independently derived
finite identities.  No novelty priority is asserted.

\paragraph{The fixed-rank Mellin object and its exact support.}
For normalized Haar measure on $U(N)$, retain
\begin{equation}
 \Phi_N(z)=\det(zI-U),\qquad
 X_N(U)=|\Phi_N'(1)|^2,\qquad
 \mathsf M_N(k)=\mathbb E X_N^k
 \label{eq:cueM-fixedN-objects}
\end{equation}
whenever the last integral is finite.  Equation
\eqref{eq:GMS-Phi-exact-phase} gives the exact convention map to
$\Lambda_N(z)=\det(1-zU^*)$: multiplication by the $z$-independent phase
$(-1)^N\overline{\det U}$ has inverse $(-1)^N\det U$, zero kernel, singleton fibres, and
preserves every derivative magnitude.  Consequently, for $s\in\mathbb Z_{>0}$,
\begin{equation}
 \mathsf M_N(s)=M_{(1^s),(1^s)}(1,N)
 \label{eq:cueM-GMS-specialization}
\end{equation}
in the notation of Grover--Mezzadri--Simm
\cite[eq.~(1.1)]{groverMezzadriSimm2026}.

\begin{theorem}[Exact fixed-rank support and Mellin representation]
\label{thm:cueM-fixedN-support}
For every $N\geq2$, the pushforward probability
\[
 \nu_N=(X_N)_*\operatorname{Haar}_{U(N)}
\]
has exact support
\begin{equation}
 \operatorname{supp}\nu_N=[0,L_N],\qquad L_N=N^2 4^{N-1}.
 \label{eq:cueM-exact-support}
\end{equation}
For every $k\geq0$,
\begin{equation}
 \mathsf M_N(k)=\int_0^{L_N}t^k\,d\nu_N(t).
 \label{eq:cueM-Mellin-measure}
\end{equation}
The same identity is tautological for any other real $k$ for which the integral is finite;
the theorem makes no additional negative-moment assertion.
\end{theorem}

\begin{proof}
If $z_1,\ldots,z_N\in\mathbb T$ are the eigenvalues, then
\[
 \Phi_N'(1)=\sum_{m=1}^{N}\prod_{j\ne m}(1-z_j).
\]
Every factor has modulus at most $2$, so
$|\Phi_N'(1)|\leq N2^{N-1}$ and $0\leq X_N\leq L_N$.  At $U=I_N$,
$\Phi_N(z)=(z-1)^N$ and $X_N=0$; at $U=-I_N$,
$\Phi_N(z)=(z+1)^N$ and $X_N=L_N$.  The continuous image of the compact connected group
$U(N)$ in $\mathbb R$ is a compact interval, hence it is exactly $[0,L_N]$.

It remains to distinguish image from measure support.  If $y=X_N(U_0)$ and $V$ is any open
interval containing $y$, continuity makes $X_N^{-1}(V)$ an open neighbourhood of $U_0$.
Normalized Haar measure has full support, so this neighbourhood has positive measure.
Thus every image point belongs to the pushforward support.  The reverse inclusion was already
proved by the bound.  Equation \eqref{eq:cueM-Mellin-measure} is the definition of a
pushforward integral.
\end{proof}

\paragraph{The complete moment cone and a constructive inverse.}
Push $\nu_N$ forward by $t\mapsto t/L_N$ and write the resulting probability as
$\widetilde\nu_N$.  Its integer moments are
\begin{equation}
 c_{N,n}=\int_0^1x^n\,d\widetilde\nu_N(x)
        =\frac{\mathsf M_N(n)}{L_N^n}.
 \label{eq:cueM-normalized-moments}
\end{equation}
Use the forward difference
$\Delta c_{N,n}=c_{N,n+1}-c_{N,n}$ without changing this sign convention.

\begin{theorem}[Finite differences, positivity, and constructive recovery]
\label{thm:cueM-Hausdorff-Bernstein}
For all $n,r\geq0$,
\begin{equation}
 (-1)^r\Delta^r c_{N,n}
 =\int_0^1x^n(1-x)^r\,d\widetilde\nu_N(x)\geq0.
 \label{eq:cueM-Hausdorff-differences}
\end{equation}
For $m\geq0$ and $0\leq j\leq m$, define
\begin{align}
 \lambda_{N,m,j}
 &=\binom mj(-1)^{m-j}\Delta^{m-j}c_{N,j} \notag\\
 &=\binom mj\int_0^1x^j(1-x)^{m-j}\,d\widetilde\nu_N(x),
 \label{eq:cueM-Bernstein-weights}\\
 \widetilde\nu_{N,m}&=\sum_{j=0}^{m}\lambda_{N,m,j}\delta_{j/m}.
 \label{eq:cueM-Bernstein-discrete}
\end{align}
Then the weights are nonnegative and sum to one, and
\begin{equation}
 \widetilde\nu_{N,m}\Rightarrow\widetilde\nu_N.
 \label{eq:cueM-Bernstein-limit}
\end{equation}
Equivalently, the integer sequence in \eqref{eq:cueM-normalized-moments} constructively
recovers the complete compact measure and therefore every nonnegative real moment:
\begin{equation}
 \frac{\mathsf M_N(k)}{L_N^k}
 =\lim_{m\to\infty}\sum_{j=0}^{m}
   \left(\frac jm\right)^k\lambda_{N,m,j},
 \qquad k\geq0.
 \label{eq:cueM-fractional-recovery}
\end{equation}
\end{theorem}

\begin{proof}
Expanding the difference gives
\[
 (-1)^r\Delta^r c_{N,n}
 =\sum_{q=0}^{r}\binom rq(-1)^q c_{N,n+q}
 =\int_0^1x^n\sum_{q=0}^{r}\binom rq(-x)^q\,d\widetilde\nu_N(x),
\]
and the binomial sum is $(1-x)^r$.  This proves
\eqref{eq:cueM-Hausdorff-differences} and both expressions in
\eqref{eq:cueM-Bernstein-weights}.  The weights are therefore nonnegative, while
\[
 \sum_{j=0}^{m}\lambda_{N,m,j}
 =\int_0^1\sum_{j=0}^{m}\binom mjx^j(1-x)^{m-j}\,d\widetilde\nu_N(x)=1.
\]
For $f\in C([0,1])$,
\[
 \int f\,d\widetilde\nu_{N,m}
 =\int_0^1\sum_{j=0}^{m}f(j/m)\binom mjx^j(1-x)^{m-j}
   \,d\widetilde\nu_N(x).
\]
The inner sum is Bernstein's polynomial $B_mf$, and
$\lVert B_mf-f\rVert_\infty\to0$ by Bernstein's theorem
\cite[eqs.~(1)--(5)]{bernstein1912}.  This proves weak convergence.  Finally take
$f(x)=x^k$, which is continuous for $k\geq0$.  This is the specialization of the general
coordinate-preserving reconstruction in
\cref{prop:bernstein-measure-reconstruction}, now written explicitly in the CUE moments.
\end{proof}

Two positivity consequences retain useful information that the scalar moments alone can
hide.  For every $(a_0,\ldots,a_m)\in\mathbb C^{m+1}$,
\begin{equation}
 \sum_{i,j=0}^{m}a_i\overline{a_j}\mathsf M_N(i+j)
 =\int_0^{L_N}\left|\sum_{i=0}^{m}a_it^i\right|^2d\nu_N(t)\geq0,
 \label{eq:cueM-Hankel-PSD}
\end{equation}
so every moment Hankel matrix is positive semidefinite.  H\"older's inequality gives
\begin{equation}
 \mathsf M_N(\theta k_1+(1-\theta)k_0)
 \leq\mathsf M_N(k_1)^\theta\mathsf M_N(k_0)^{1-\theta}
 \label{eq:cueM-log-convex}
\end{equation}
whenever both endpoint moments are finite and $0\leq\theta\leq1$.

\begin{proposition}[Stieltjes generating function with exact radius]
\label{prop:cueM-Stieltjes-radius}
The series
\[
 G_N(\tau)=\sum_{n\geq0}\mathsf M_N(n)\tau^n
\]
has radius $L_N^{-1}$ and, for $|\tau|<L_N^{-1}$,
\begin{equation}
 G_N(\tau)=\int_0^{L_N}\frac{d\nu_N(t)}{1-\tau t}.
 \label{eq:cueM-Stieltjes-transform}
\end{equation}
\end{proposition}

\begin{proof}
The upper support bound gives $\mathsf M_N(n)\leq L_N^n$.  Conversely, for every
$\varepsilon>0$, exact support gives
$\nu_N([L_N-\varepsilon,L_N])>0$ and hence
\[
 \mathsf M_N(n)\geq
 \nu_N([L_N-\varepsilon,L_N])(L_N-\varepsilon)^n.
\]
Therefore $\lim_n\mathsf M_N(n)^{1/n}=L_N$.  Cauchy--Hadamard proves the radius, and uniform
convergence of the geometric series on $[0,L_N]$ proves
\eqref{eq:cueM-Stieltjes-transform}.
\end{proof}

\paragraph{Several large-rank moment scales.}
For each fixed positive integer $s$, Grover--Mezzadri--Simm prove
\begin{equation}
 \mathsf M_N(s)\sim b_sN^{s^2+2s},\qquad b_s>0
 \label{eq:cueM-GMS-exponent}
\end{equation}
as the $(1^s),(1^s)$ instance of their microscopic theorem
\cite[Thm.~1.2 and Sec.~2]{groverMezzadriSimm2026}; the complete source-level determinant
audit is in \cref{sec:GMS-complete-proof-audit}.

\begin{theorem}[No deterministic scale carries two convergent nonzero moments]
\label{thm:cueM-no-one-scale}
If $r,s\in\mathbb Z_{>0}$ and $r\ne s$, there are no $a_N>0$ and
$c_r,c_s\in(0,\infty)$ such that
\begin{equation}
 \frac{\mathsf M_N(r)}{a_N^r}\longrightarrow c_r,
 \qquad
 \frac{\mathsf M_N(s)}{a_N^s}\longrightarrow c_s.
 \label{eq:cueM-two-moment-scale}
\end{equation}
For every fixed $m\geq1$ and $\alpha>0$,
\begin{equation}
 \mathbb P(X_N\geq N^\alpha)
 \leq N^{m^2+2m-\alpha m+o(1)}.
 \label{eq:cueM-Markov-tail}
\end{equation}
For $0<\theta<1$,
\begin{equation}
 \mathbb P\!\left(
 X_N\geq(\theta b_m)^{1/m}N^{m+2}(1+o(1))
 \right)
 \geq(1-\theta)^2\frac{b_m^2}{b_{2m}}N^{-2m^2+o(1)}.
 \label{eq:cueM-PZ-tail}
\end{equation}
\end{theorem}

\begin{proof}
If \eqref{eq:cueM-two-moment-scale} held, then
$a_N^r\sim(b_r/c_r)N^{r^2+2r}$ and
$a_N^s\sim(b_s/c_s)N^{s^2+2s}$.  Taking logarithms and dividing by $\log N$ would make the
same sequence $\log a_N/\log N$ converge to both $r+2$ and $s+2$.

Markov's inequality applied to $X_N^m$ gives
\[
 \mathbb P(X_N\geq N^\alpha)
 \leq\mathsf M_N(m)N^{-\alpha m},
\]
which becomes \eqref{eq:cueM-Markov-tail}.  For completeness, the inequality used in the
other direction is obtained directly.  If $Z\geq0$, then
\[
 \mathbb EZ\leq\theta\mathbb EZ+
 \mathbb E[Z\mathbf1_{\{Z\geq\theta\mathbb EZ\}}]
 \leq\theta\mathbb EZ+(\mathbb EZ^2)^{1/2}
 \mathbb P(Z\geq\theta\mathbb EZ)^{1/2}.
\]
After rearrangement this is the Paley--Zygmund bound
\cite{paleyZygmund1932}.  Put $Z=X_N^m$ and use
\eqref{eq:cueM-GMS-exponent} at $m$ and $2m$ to obtain
\eqref{eq:cueM-PZ-tail}.
\end{proof}

The maximal-support coordinate illustrates the same multiscale behaviour without an
asymptotic theorem.  Schur-character orthogonality gives
\begin{equation}
 \mathsf M_N(1)=\sum_{j=0}^{N-1}(N-j)^2
 =\frac{N(N+1)(2N+1)}6.
 \label{eq:cueM-first-moment}
\end{equation}
Indeed,
$\Phi_N(z)=\sum_{j=0}^{N}(-1)^je_j(U)z^{N-j}$ and the exterior-power characters satisfy
$\mathbb E[e_j(U)\overline{e_\ell(U)}]=\delta_{j\ell}$; this is the same Schur-coordinate
orthogonality used throughout the derivative-moment literature
\cite{macdonald1995,dehaye2008,groverMezzadriSimm2026}.  Hence
$\mathbb E(X_N/L_N)\to0$.  Since $0\leq X_N/L_N\leq1$, Markov's inequality proves
$X_N/L_N\to0$ in probability and therefore weakly.  This conclusion concerns the retained
maximal-support scale $L_N$; it does not replace the distinct microscopic regimes in
Grover--Mezzadri--Simm or Simm--Wei.

\paragraph{The measure-valued Verblunsky recursion.}
Define the state domain
\begin{equation}
 \mathcal A_j=\{u\in\mathbb C:1-j+2\operatorname{Re}u>0\},
 \qquad 0\leq j\leq N-1.
 \label{eq:cueM-state-domain}
\end{equation}
It contains $u_0=0$.  For $u\in\mathcal A_j$ retain, without changing variables,
\begin{align}
 T_j(\beta;u)&=\frac{1+u-\beta(j-\overline u)}{1-\beta},
 \label{eq:cueM-T-state}\\
 \eta_j(u)&=\frac{N-1-\overline u}{N-j+u},
 \label{eq:cueM-eta-state}\\
 \gamma_j(\beta,u)&=|1-\beta\eta_j(u)|^2.
 \label{eq:cueM-gamma-state}
\end{align}
For $0\leq j\leq N-2$, let
\begin{equation}
 d\mu_{N-j-1}(\beta)=
 \frac{N-j-1}{\pi}(1-|\beta|^2)^{N-j-2}\mathbf1_{\mathbb D}(\beta)d^2\beta,
 \label{eq:cueM-beta-law}
\end{equation}
and take the terminal $\beta_{N-1}$ uniformly on $\mathbb T$.  These are exactly the
independent pushforward laws proved in
\cref{thm:verblunsky-peel-derived}, not fitted approximations.

\begin{theorem}[State kernel, boundary law, and measure peel]
\label{thm:cueM-measure-peel}
The map $T_j(\beta;\cdot)$ sends $\mathcal A_j$ into $\mathcal A_{j+1}$ for every
$\beta\in\mathbb D$, and $|\eta_j(u)|<1$ on $\mathcal A_j$.  Starting at a state
$u\in\mathcal A_j$, define the future product recursively by
\begin{align}
 Y_{N-1}(u)&=|1-e^{-\mathrm i\Theta}\eta_{N-1}(u)|^2,
 \label{eq:cueM-terminal-Y}\\
 Y_j(u)&=\gamma_j(\beta_j,u)
 Y_{j+1}(T_j(\beta_j;u)),
 \qquad j\leq N-2,
 \label{eq:cueM-recursive-Y}
\end{align}
where $\Theta$ is uniform on $[0,2\pi]$ and all future variables are independent.  Let
$\mu_{j,u}$ be the law of $Y_j(u)$.  Then
\begin{equation}
 \mu_{j,u}=
 \int_{\mathbb D}\sigma_{\gamma_j(\beta,u)}
 \mu_{j+1,T_j(\beta;u)}\,d\mu_{N-j-1}(\beta),
 \label{eq:cueM-measure-recursion}
\end{equation}
where $\sigma_a\mu=(x\mapsto ax)_*\mu$.  This equality means equality against every bounded
continuous test function and therefore specifies the map, its domain, and its codomain.
Moreover,
\begin{equation}
 \mathsf M_N(k)=N^{2k}\int y^k\,d\mu_{0,0}(y),
 \qquad k\geq0.
 \label{eq:cueM-peel-moment}
\end{equation}

At the boundary put $q=|\eta_{N-1}(u)|$.  If $q>0$, $\mu_{N-1,u}$ has density
\begin{equation}
 \frac{
 \mathbf1_{[(1-q)^2,(1+q)^2]}(y)
 }{\pi\sqrt{(y-(1-q)^2)((1+q)^2-y)}}\,dy;
 \label{eq:cueM-boundary-arcsine}
\end{equation}
if $q=0$, it is the atom $\delta_1$.  For every $k\geq0$,
\begin{equation}
 \int y^k\,d\mu_{N-1,u}(y)
 ={}_2F_1(-k,-k;1;q^2).
 \label{eq:cueM-boundary-hypergeometric}
\end{equation}
\end{theorem}

\begin{proof}
Put $w=2u-j$.  Direct substitution in \eqref{eq:cueM-T-state} gives
\begin{equation}
 1+\operatorname{Re}(2T_j(\beta;u)-(j+1))
 =\frac{1-|\beta|^2}{|1-\beta|^2}
  (1+\operatorname{Re}w)>0,
 \label{eq:cueM-state-propagation}
\end{equation}
so the state domain is preserved.  A second direct calculation gives
\begin{equation}
 |N-j+u|^2-|N-1-\overline u|^2
 =(2N-j-1)(1-j+2\operatorname{Re}u)>0,
 \label{eq:cueM-eta-bound}
\end{equation}
which proves both nonvanishing of the denominator and $|\eta_j(u)|<1$.

Conditioning \eqref{eq:cueM-recursive-Y} on $\beta_j=\beta$ gives, for bounded continuous
$f$,
\[
 \mathbb E[f(Y_j(u))\mid\beta_j=\beta]
 =\int f(\gamma_j(\beta,u)y)\,
 d\mu_{j+1,T_j(\beta;u)}(y).
\]
Integration in $\beta$ proves \eqref{eq:cueM-measure-recursion}.  The pathwise product in
\eqref{eq:KN-complete-peel-modulus} proves \eqref{eq:cueM-peel-moment}; no moment recursion
has been substituted for the underlying random variable.

At the boundary,
$Y=1+q^2-2q\cos\Theta$.  The pushforward of uniform angle under $x=\cos\Theta$ has density
$(\pi\sqrt{1-x^2})^{-1}$ on $(-1,1)$.  The affine map
$y=1+q^2-2qx$ has inverse $x=(1+q^2-y)/(2q)$ and derivative magnitude $(2q)^{-1}$;
substitution gives \eqref{eq:cueM-boundary-arcsine}.  If $q=0$, the map is constantly one.
Finally, $q<1$ and the two binomial series converge absolutely:
\[
 \frac1{2\pi}\int_0^{2\pi}
 (1-qe^{-\mathrm i\theta})^k(1-qe^{\mathrm i\theta})^k\,d\theta
 =\sum_{n\geq0}\frac{(-k)_n^2}{(n!)^2}q^{2n},
\]
because angular integration retains exactly the equal Fourier indices.  This is
\eqref{eq:cueM-boundary-hypergeometric}; $q=0$ is included by direct evaluation.
\end{proof}

\paragraph{The exact transform carrier.}
Let $\mathcal C_c$ be the real vector space of finite signed compactly supported Borel
measures on $[0,\infty)$, with additive convolution $*$ and unit $\delta_0$.  For
$\mu\in\mathcal C_c$, $t\geq0$, and $\operatorname{Re}s>1$, define
\begin{align}
 \mathcal L[\mu](t)&=\int e^{-tx}\,d\mu(x),
 \label{eq:cueM-Laplace}\\
 \mathcal Z_s[\mu](t)&=\int\zeta(s,1+tx)\,d\mu(x),
 \label{eq:cueM-Hurwitz}\\
 \mathcal D[\mu](t)&=\int\psi(1+tx)\,d\mu(x).
 \label{eq:cueM-digamma}
\end{align}

\begin{theorem}[Transported convolution algebra and Bose coordinate]
\label{thm:cueM-transform-algebra}
The Laplace map is an injective algebra homomorphism
\begin{equation}
 \mathcal L[\mu*\nu]=\mathcal L[\mu]\mathcal L[\nu].
 \label{eq:cueM-Laplace-product}
\end{equation}
For
\begin{equation}
 (\mathcal B_sF)(t)=\frac1{\Gamma(s)}
 \int_0^\infty\frac{x^{s-1}}{e^x-1}F(tx)\,dx,
 \label{eq:cueM-Bose-operator}
\end{equation}
one has
\begin{equation}
 \mathcal B_s\mathcal L[\mu]=\mathcal Z_s[\mu],
 \qquad
 -\operatorname{FP}_{s=1}\mathcal Z_s[\mu]=\mathcal D[\mu].
 \label{eq:cueM-Bose-Hurwitz-digamma}
\end{equation}
Put
\[
 \mathfrak C_s=\{(\mathcal L[\mu],\mathcal Z_s[\mu],\mathcal D[\mu]):
 \mu\in\mathcal C_c\}.
\]
Then
\begin{equation}
 \mathbf T_s:\mathcal C_c\longrightarrow
 \mathfrak C_s,\qquad
 \mu\longmapsto(\mathcal L[\mu],\mathcal Z_s[\mu],\mathcal D[\mu]),
 \label{eq:cueM-transform-triple}
\end{equation}
is a bijection onto its displayed image, with zero kernel and singleton fibres.  Transporting
$*$ through this bijection makes $\mathfrak C_s$ a commutative unital algebra.

For $a\geq0$ let $S_aF(t)=F(at)$.  Then
\begin{equation}
 \mathbf T_s[\sigma_a\mu]=S_a\mathbf T_s[\mu],
 \qquad
 \mathcal B_sS_a=S_a\mathcal B_s,
 \label{eq:cueM-dilation-commutation}
\end{equation}
and the measure peel becomes the exact linear recursion
\begin{equation}
 \mathbf T_s[\mu_{j,u}]=
 \int_{\mathbb D}S_{\gamma_j(\beta,u)}
 \mathbf T_s[\mu_{j+1,T_j(\beta;u)}]\,
 d\mu_{N-j-1}(\beta).
 \label{eq:cueM-transform-peel}
\end{equation}
This is an averaging of dilation operators on a transported convolution algebra; it is not
being identified with convolution multiplication.
\end{theorem}

\begin{proof}
Fubini gives
\[
 \mathcal L[\mu*\nu](t)=\iint e^{-t(x+y)}\,d\mu(x)d\nu(y)
 =\mathcal L[\mu](t)\mathcal L[\nu](t).
\]
If $\mathcal L[\mu]$ vanishes identically and $\operatorname{supp}\mu\subset[0,R]$, then
all derivatives at zero vanish, so $\int x^n\,d\mu(x)=0$ for every $n$.  Polynomial density
in $C([0,R])$ and the Riesz representation theorem imply $\mu=0$; equivalently this is the
compact-support instance of Fourier--Stieltjes uniqueness
\cite{rudin1990}.  Thus the first component of \eqref{eq:cueM-transform-triple} already
recovers $\mu$.

 The classical Mellin integral
\begin{equation}
 \zeta(s,a)=\frac1{\Gamma(s)}\int_0^\infty
 \frac{x^{s-1}e^{-(a-1)x}}{e^x-1}\,dx,
 \qquad \operatorname{Re}s>1, a>0,
 \label{eq:cueM-Hurwitz-Mellin}
\end{equation}
 and the simple pole at $s=1$ with residue one are recorded in the canonical Hurwitz source
 \cite[\S25.11(i),(vii)]{apostolDLMF25}.  The constant term needed here can be derived
 without importing an unstated formula.  For $a>0$,
 \[
  D(s,a):=\zeta(s,a)-\zeta(s)
  =\sum_{n=0}^{\infty}\bigl((n+a)^{-s}-(n+1)^{-s}\bigr).
 \]
 The series and all of its $s$-derivatives converge locally uniformly for
 $\operatorname{Re}s>0$: on a compact $a$-interval the mean-value theorem bounds the
 $r$th derivative of the $n$th difference by
 $C_r(1+\log^{r+1}(n+1))(n+1)^{-\operatorname{Re}s-1}$.  Hence $D$ is holomorphic near
 $s=1$.  Using the Laurent expansion
 $\zeta(s)=(s-1)^{-1}+\gamma+O(s-1)$ and the digamma series
 \cite[eq.~(5.7.6)]{askeyRoyDLMF5},
 \[
  D(1,a)=\sum_{n=0}^{\infty}\left(\frac1{n+a}-\frac1{n+1}\right)
        =-\psi(a)-\gamma,
 \]
 so indeed $\zeta(s,a)=(s-1)^{-1}-\psi(a)+O(s-1)$, locally uniformly for $a>0$.
 Substitute $a=1+ty$ in
\eqref{eq:cueM-Hurwitz-Mellin}, integrate in $y$, and use compact support to justify Fubini;
this proves the first identity in \eqref{eq:cueM-Bose-Hurwitz-digamma}.  Integrating the
Laurent expansion uniformly for $a$ in the resulting compact interval gives
\[
 \mathcal Z_s[\mu](t)=\frac{\mu([0,\infty))}{s-1}
 -\mathcal D[\mu](t)+O(s-1),
\]
which proves the finite-part identity.  Bijectivity onto the displayed image and the
transported product now follow from injectivity of $\mathcal L$.

Finally, every transform in \eqref{eq:cueM-Laplace}--\eqref{eq:cueM-digamma} satisfies
$F_{\sigma_a\mu}(t)=F_\mu(at)$.  This proves the first dilation equation; direct substitution
in \eqref{eq:cueM-Bose-operator} proves the second.  Apply these identities to
\eqref{eq:cueM-measure-recursion} to obtain \eqref{eq:cueM-transform-peel}.
\end{proof}

\paragraph{The complete noninteger $U(2)$ moment.}
The exact Haar coordinate theorem \cref{thm:u2-exact-derivative-law} gives independent
\begin{equation}
 R\sim\operatorname{Beta}\!\left(\frac12,\frac32\right),
 \qquad \phi\sim\operatorname{Unif}(\mathbb R/2\pi\mathbb Z),
 \qquad X_2=4|1-\sqrt R e^{\mathrm i\phi}|^2,
 \label{eq:cueM-U2-law}
\end{equation}
including the two-sheet folding map, its boundary fibres, and the exact support $[0,16]$.

\begin{corollary}[Closed $U(2)$ Mellin moment]
\label{cor:cueM-U2-3F2}
For every real $k\geq0$,
\begin{equation}
 \mathsf M_2(k)=4^k\,{}_3F_2\!\left(
 \begin{matrix}-k,-k,\frac12\\1,2\end{matrix};1\right).
 \label{eq:cueM-U2-3F2}
\end{equation}
\end{corollary}

\begin{proof}
Conditioning on $R=r$ and integrating the uniform phase gives
\begin{equation}
 \frac1{2\pi}\int_0^{2\pi}|1-\sqrt r e^{\mathrm i\phi}|^{2k}\,d\phi
 ={}_2F_1(-k,-k;1;r).
 \label{eq:cueM-U2-angular}
\end{equation}
Indeed, expand both binomial factors; the angular integral deletes precisely the unequal
indices.  The beta density in \eqref{eq:u2-beta-density} therefore yields
\begin{equation}
 \mathsf M_2(k)=\frac{2^{2k+1}}\pi
 \int_0^1r^{-1/2}(1-r)^{1/2}{}_2F_1(-k,-k;1;r)\,dr.
 \label{eq:cueM-U2-beta-integral}
\end{equation}
For $k\geq0$ the hypergeometric coefficients are nonnegative, so Tonelli applies.  Termwise
beta integration gives
\[
 \int_0^1r^{n-1/2}(1-r)^{1/2}\,dr
 =B\!\left(n+\frac12,\frac32\right)
 =\frac\pi2\frac{(1/2)_n}{(2)_n}.
\]
Substitution in \eqref{eq:cueM-U2-beta-integral} is exactly the defining ${}_3F_2$ series in
\eqref{eq:cueM-U2-3F2}; the hypergeometric conventions and the terminating
Chu--Vandermonde identity are recorded in the canonical reference
\cite{oldeDaalhuisDLMF15}.  Independently, the finite certificate below compares the
terminating right-hand side with Weyl constant terms for $k=0,\ldots,7$, obtaining
\[
 1,5,34,285,2766,29666,340740,4108509.
\]
\end{proof}

\paragraph{The first peeled $U(3)$ compact law.}
Let $\beta_0\in\mathbb D$ and retain the exact rational coordinates
\begin{equation}
 \lambda=
 \frac{(1-2\overline{\beta_0})(1-\beta_0)}
 {(3-2\beta_0)(1-\overline{\beta_0})},\qquad
 c=\frac{\overline{\beta_0}(1-\beta_0)}
 {(3-2\beta_0)(1-\overline{\beta_0})},\qquad
 d=\frac{2-\beta_0}{3-2\beta_0}.
 \label{eq:cueM-U3-lambda-c-d}
\end{equation}
These are the same coordinates as \eqref{eq:n3-peeled-lambda-c-d}; their base quotient,
inverse, fibres, branch locus, and physical $J$-sheet are proved without omissions in
\cref{thm:n3-base-quotient,thm:n3-physical-trace-cover}.

The connection to the state recursion is also exact.  Starting with $u_0=0$ gives
\begin{equation}
 u_1=T_0(\beta_0;0)=\frac1{1-\beta_0},\qquad
 \eta_1(u_1)=\lambda.
 \label{eq:cueM-U3-first-state}
\end{equation}
For the next disk variable $\beta\in\mathbb D$, put
$u_2=T_1(\beta;u_1)$.  Substitution in \eqref{eq:cueM-T-state} and
\eqref{eq:cueM-eta-state} gives the retained phase relation
\begin{equation}
 \eta_2(u_2)=
 -\frac{\beta-1}{\overline\beta-1}
 \frac{c+\overline\beta d}{1-\lambda\beta}.
 \label{eq:cueM-U3-second-state}
\end{equation}
The first factor on the right has modulus one.  Hence, with
\begin{equation}
 A_{\beta_0}(\beta)=|1-\lambda\beta|^2,\qquad
 B_{\beta_0}(\beta)=|c+\overline\beta d|^2,
 \label{eq:cueM-U3-A-B}
\end{equation}
one has $B/A=|\eta_2(u_2)|^2<1$.  This proves the precise state-space arrow behind the local
inner kernel; it is not an analogy between two formulas.

\begin{theorem}[Local and global $U(3)$ compact measures]
\label{thm:cueM-U3-compact-measures}
For $\phi\in[0,\pi]$ define
\begin{equation}
 \Theta_{\beta_0}(\beta,\phi)=
 A_{\beta_0}(\beta)+B_{\beta_0}(\beta)
 +2\sqrt{A_{\beta_0}(\beta)B_{\beta_0}(\beta)}\cos\phi.
 \label{eq:cueM-U3-Theta}
\end{equation}
Then, for $k\geq0$, the local inner kernel is the moment
\begin{align}
 \Psi_k(\beta_0)
 &=\frac1\pi\int_{\mathbb D}
 A_{\beta_0}(\beta)^k
 {}_2F_1\!\left(-k,-k;1;\frac{B_{\beta_0}(\beta)}
 {A_{\beta_0}(\beta)}\right)d^2\beta
 \label{eq:cueM-U3-Psi-hypergeometric}\\
 &=\frac1{\pi^2}\int_{\mathbb D}\int_0^\pi
 \Theta_{\beta_0}(\beta,\phi)^k\,d\phi\,d^2\beta.
 \label{eq:cueM-U3-Psi-measure}
\end{align}
For fixed $\beta_0,\beta$, if $B>0$ the conditional pushforward in $t$ has density
\begin{equation}
 \frac{
 \mathbf1_{[(\sqrt A-\sqrt B)^2,(\sqrt A+\sqrt B)^2]}(t)
 }{\pi\sqrt{(t-(\sqrt A-\sqrt B)^2)
 ((\sqrt A+\sqrt B)^2-t)}}\,dt;
 \label{eq:cueM-U3-conditional-density}
\end{equation}
if $B=0$ it is $\delta_A$.

Define
\begin{equation}
 \Xi(\beta_0,\beta,\phi)=|3-2\beta_0|^2
 \Theta_{\beta_0}(\beta,\phi)
 \label{eq:cueM-U3-Xi}
\end{equation}
and push the probability measure
\begin{equation}
 \frac2\pi(1-|\beta_0|^2)d^2\beta_0\otimes
 \frac{d^2\beta}{\pi}\otimes\frac{d\phi}{\pi}
 \label{eq:cueM-U3-product-measure}
\end{equation}
forward under $\Xi$.  The result is exactly $\nu_3$ from
\cref{thm:cueM-fixedN-support}; in particular
\begin{equation}
 \mathsf M_3(k)=\int_0^{144}t^k\,d\nu_3(t),\qquad
 \operatorname{supp}\nu_3=[0,144].
 \label{eq:cueM-U3-full-measure}
\end{equation}
\end{theorem}

\begin{proof}
Fix $\beta_0,\beta$ and abbreviate $A=A_{\beta_0}(\beta)>0$ and
$r=\sqrt{B/A}<1$.  Replacing a uniform angle by its translate if necessary,
\[
 A+B+2\sqrt{AB}\cos\phi=A|1+re^{\mathrm i\phi}|^2.
\]
The same diagonal Fourier calculation as in
\eqref{eq:cueM-U2-angular} gives
\[
 \frac1\pi\int_0^\pi
 (A+B+2\sqrt{AB}\cos\phi)^k\,d\phi
 =A^k{}_2F_1(-k,-k;1;B/A),
\]
proving \eqref{eq:cueM-U3-Psi-measure}.  The change of variables from
$x=\cos\phi$ to $t=A+B+2\sqrt{AB}x$ gives
\eqref{eq:cueM-U3-conditional-density}; when $B=0$ the map is constant.

For the outer step, the first factor of the pathwise derivative peel is
$|1-(2/3)\beta_0|^2$.  Multiplying by the retained $3^2$ in
\eqref{eq:cueM-peel-moment} gives $|3-2\beta_0|^2$.  The next disk law is uniform on
$\mathbb D$, and the last circle law is precisely the angle just integrated.  Therefore the
pushforward of \eqref{eq:cueM-U3-product-measure} is the law of $X_3$, proving the first
assertion in \eqref{eq:cueM-U3-full-measure}.  Its exact support now follows from
\cref{thm:cueM-fixedN-support} with $L_3=9\cdot16=144$; the local source's direct inequality
$\Xi\leq144$ is consistent with, but weaker than, this equality.
\end{proof}

\paragraph{Integer kernels on the retained phase cover.}
Put
\begin{equation}
 u=|\lambda|^2,\qquad v=|c|^2,\qquad
 w=|d|^2=\frac{1-u+2v}{2},\qquad
 Z=c\overline d\,\overline\lambda.
 \label{eq:cueM-U3-u-v-w-Z}
\end{equation}
The complete $(u,v,J)$ cover, including the inverse to $\beta_0$, is already proved in
\cref{thm:n3-physical-trace-cover}; here the retained $Z$ coordinate obeys
\begin{equation}
 Z+\overline Z=R(u,v)=\frac{9u^2-16uv-10u-16v+1}{16},\qquad
 Z\overline Z=P(u,v)=uvw.
 \label{eq:cueM-U3-trace-norm}
\end{equation}
Choose $a=\sqrt u$, $b=\sqrt v$, and $\rho=\sqrt w$.  If $ab\rho>0$, choose
$\delta$ modulo $2\pi$ by $Z=ab\rho e^{-\mathrm i\delta}$; if $ab\rho=0$, $Z=0$ and
$\delta$ is an immaterial fibre coordinate.  For $s\in\mathbb Z_{\geq0}$ and
$0\leq m\leq s$, set
\begin{align}
 P_{s,m}(z)&=(1-az)^{s-m}(b+\rho e^{-\mathrm i\delta}z)^m
 =\sum_{n=0}^{s}p_{s,m,n}z^n,
 \label{eq:cueM-U3-Psm}\\
 p_{s,m,n}&=\sum_{j=0}^{\min(m,n)}
 (-a)^{n-j}b^{m-j}\rho^je^{-\mathrm i j\delta}
 \binom{s-m}{n-j}\binom mj.
 \label{eq:cueM-U3-psmn}
\end{align}

\begin{theorem}[Hermitian square and complete trace descent]
\label{thm:cueM-U3-trace-descent}
For every integer $s\geq0$,
\begin{equation}
 \Psi_s=\sum_{m=0}^{s}\binom sm^2
 \sum_{n=0}^{s}\frac{|p_{s,m,n}|^2}{n+1}.
 \label{eq:cueM-U3-Hermitian-square}
\end{equation}
Define
\begin{equation}
 T_0=2,\qquad T_1=R,\qquad T_{\ell+1}=RT_\ell-PT_{\ell-1}.
 \label{eq:cueM-U3-trace-recurrence}
\end{equation}
Then $T_\ell=Z^\ell+\overline Z^{\,\ell}$ and the phase-invariant kernel is the explicit
polynomial
\begin{align}
 \Psi_s^\sharp(u,v)
 &=\sum_{m=0}^{s}\binom sm^2\left(
 A_{s,m}+\sum_{\ell=1}^{\min(m,s-m)}(-1)^\ell T_\ell B_{s,m,\ell}
 \right),
 \label{eq:cueM-U3-Psi-general}\\
 A_{s,m}
 &=\sum_{p=0}^{s-m}\sum_{q=0}^{m}
 \frac{\binom{s-m}{p}^2\binom mq^2}{p+q+1}
 u^pv^{m-q}w^q,
 \label{eq:cueM-U3-A-sm}\\
 B_{s,m,\ell}
 &=\sum_{p=0}^{s-m-\ell}\sum_{q=0}^{m-\ell}
 \frac{
 \binom{s-m}{p+\ell}\binom{s-m}{p}
 \binom m{q+\ell}\binom mq
 }{p+q+\ell+1}
 u^pv^{m-q-\ell}w^q.
 \label{eq:cueM-U3-B-sml}
\end{align}
It belongs to $\mathbb Q[u,v]$, has total degree exactly $s$, and
\begin{equation}
 [v^s]\Psi_s^\sharp(u,v)=\frac1{s+1}\binom{2s+1}{s}.
 \label{eq:cueM-U3-v-leading}
\end{equation}
\end{theorem}

\begin{proof}
For integer $s$,
${}_2F_1(-s,-s;1;x)=\sum_{m=0}^{s}\binom sm^2x^m$.  Inserting this terminating identity
in \eqref{eq:cueM-U3-Psi-hypergeometric} turns the $m$th summand into
\[
 |1-a\beta|^{2(s-m)}|b+\rho e^{\mathrm i\delta}\overline\beta|^{2m}
 =|P_{s,m}(\beta)|^2.
\]
The exact disk orthogonality
\[
 \frac1\pi\int_{\mathbb D}\beta^n\overline\beta^{\,r}d^2\beta
 =\frac{\delta_{nr}}{n+1}
\]
proves \eqref{eq:cueM-U3-Hermitian-square} after substituting
\eqref{eq:cueM-U3-psmn}.

The quadratic equation
$Z^2-RZ+P=0$ and its conjugate imply
$Z^{\ell+1}+\overline Z^{\,\ell+1}
=R(Z^\ell+\overline Z^{\,\ell})-P(Z^{\ell-1}+\overline Z^{\,\ell-1})$;
induction proves the trace recurrence.  Now expand $|p_{s,m,n}|^2$.  Terms with equal phase
indices give \eqref{eq:cueM-U3-A-sm}.  If the two indices differ by $\ell>0$, the term and its
conjugate combine to
$(-1)^\ell(Z^\ell+\overline Z^{\,\ell})$ times the coefficient in
\eqref{eq:cueM-U3-B-sml}.  Summing all allowed indices proves
\eqref{eq:cueM-U3-Psi-general}; because $w,R,P\in\mathbb Q[u,v]$, the recurrence puts every
$T_\ell$ and therefore $\Psi_s^\sharp$ in $\mathbb Q[u,v]$.

For the degree bound, $A_{s,m}$ has degree at most $s$.  Induction in
\eqref{eq:cueM-U3-trace-recurrence} gives $\deg T_\ell\leq2\ell$, while
$\deg B_{s,m,\ell}\leq s-2\ell$; hence every off-diagonal term also has degree at most $s$.
To isolate the pure $v^s$ coefficient, set $u=0$.  Then
$w=v+1/2$, $P=0$, $R=1/16-v$, and $T_\ell=R^\ell$ for $\ell\geq1$.  At a fixed $m$, both
the diagonal and off-diagonal contributions have $v$-degree at most $m$.  Thus only $m=s$
can contribute to $v^s$, and its coefficient is
\[
 \sum_{q=0}^{s}\frac{\binom sq^2}{q+1}
 =\frac1{s+1}\sum_{q=0}^{s}\binom sq\binom{s+1}{q+1}
 =\frac1{s+1}\binom{2s+1}{s},
\]
by Vandermonde's identity.  It is nonzero, proving both
\eqref{eq:cueM-U3-v-leading} and exact degree $s$.
\end{proof}

The first four resulting polynomials, retained rather than referred to by a private source,
are
\begin{align}
 \Psi_1^\sharp={}&\frac{5+u+6v}{4},
 \label{eq:cueM-Psi1}\\
 \Psi_2^\sharp={}&\frac{47}{24}+\frac{11}{4}u+\frac{28}{3}v
 -\frac{11}{8}u^2+4uv+\frac{10}{3}v^2,
 \label{eq:cueM-Psi2}\\
 \Psi_3^\sharp={}&\frac{105}{32}+\frac{465}{32}u+\frac{75}{2}v
 -\frac{225}{32}u^2+\frac{195}{4}uv+\frac{405}{8}v^2\notag\\
 &-\frac{65}{32}u^3-\frac{15}{2}u^2v+\frac{225}{8}uv^2+\frac{35}{4}v^3,
 \label{eq:cueM-Psi3}\\
 \Psi_4^\sharp={}&\frac1{320}\bigl(
 1759u^4-16224u^3v-11812u^3+3024u^2v^2-12816u^2v-2334u^2\notag\\
 &\quad+50176uv^3+165984uv^2+117312uv+18732u
 +8064v^4+78848v^3\notag\\
 &\quad+121296v^2+40752v+1719\bigr).
 \label{eq:cueM-Psi4}
\end{align}

\begin{proposition}[Exact sum-of-squares obstruction after descent]
\label{prop:cueM-Psi2-not-SOS}
The polynomial $\Psi_2^\sharp$ is not a sum of squares in $\mathbb R[u,v]$, although its
pullback to the retained phase cover is the Hermitian square sum
\eqref{eq:cueM-U3-Hermitian-square}.
\end{proposition}

\begin{proof}
If a real polynomial of degree at most two is a sum of polynomial squares, each summand has
degree at most one: otherwise the sum of the nonzero squares of their top homogeneous parts
would be the zero polynomial.  Thus any such representation of $\Psi_2^\sharp$ would have a
positive-semidefinite affine Gram matrix.  Coefficient comparison in the fixed ordered basis
$(1,u,v)$ gives the unique symmetric matrix
\[
 Q=\begin{pmatrix}
 47/24&11/8&14/3\\
 11/8&-11/8&2\\
 14/3&2&10/3
 \end{pmatrix},
 \qquad
 \Psi_2^\sharp=(1,u,v)Q(1,u,v)^{\mathsf T}.
\]
Its leading $2\times2$ principal minor equals
\[
 \frac{47}{24}\left(-\frac{11}{8}\right)-\left(\frac{11}{8}\right)^2
 =-\frac{55}{12}<0.
\]
Hence $Q$ is not positive semidefinite.  The conclusion concerns polynomial squares in the
quotient coordinates; equation \eqref{eq:cueM-U3-Hermitian-square} records the exact map to
the different Hermitian-square presentation.
\end{proof}

The outer integral in \eqref{eq:cueM-U3-full-measure}, or independently the Weyl
constant-term formula in \cref{prop:n3-cue-exact-moments}, gives
\begin{equation}
 \mathsf M_3(1)=14,\qquad \mathsf M_3(2)=375,\qquad
 \mathsf M_3(3)=15510,\qquad \mathsf M_3(4)=847280,\qquad
 \mathsf M_3(5)=55331424.
 \label{eq:cueM-U3-finite-moments}
\end{equation}
These five values have independent exact Python/SymPy and bounded Lean certificates recorded
in \texttt{CERT-N3-CUE-20260825-0001} and \texttt{FORMAL-20260825-0010}.

\paragraph{The earlier Gauss coefficient law, with its moment domain repaired.}
For $m\in\mathbb Z_{\geq1}$ and real $k>-1/2$, put
\begin{equation}
 F_m(z)={}_2F_1(-k,-k;m;z),\qquad
 a_n^{(m,k)}=\frac{(-k)_n^2}{(m)_n n!}.
 \label{eq:cueM-Gauss-F}
\end{equation}

\begin{theorem}[Auxiliary exponent probability law]
\label{thm:cueM-Gauss-law}
The numbers
\begin{equation}
 p_n^{(m,k)}=\frac{a_n^{(m,k)}}{F_m(1)},\qquad n\geq0,
 \label{eq:cueM-Gauss-probabilities}
\end{equation}
form a probability distribution.  If $Y_{m,k}$ has this law, then for
$0\leq\varepsilon\leq1$,
\begin{equation}
 \frac{F_m(1-\varepsilon)}{F_m(1)}
 =\mathbb E(1-\varepsilon)^{Y_{m,k}}.
 \label{eq:cueM-Gauss-PGF}
\end{equation}
If $k=s\in\mathbb Z_{\geq0}$, the support is $\{0,\ldots,s\}$; if
$k\notin\mathbb Z_{\geq0}$, the support is infinite and
\begin{equation}
 p_n^{(m,k)}\sim
 \frac{\Gamma(m)}{\Gamma(-k)^2F_m(1)}n^{-(m+2k+1)}.
 \label{eq:cueM-Gauss-tail}
\end{equation}

For $s\in\mathbb Z_{>0}$ the factorial moments are
\begin{align}
 \mathbb EY_{m,s}&=\frac{s^2}{m+2s-1},
 \label{eq:cueM-Gauss-mean-integer}\\
 \mathbb E[Y_{m,s}(Y_{m,s}-1)]
 &=\frac{s^2(s-1)^2}{(m+2s-2)(m+2s-1)}.
 \label{eq:cueM-Gauss-second-integer}
\end{align}
For $s=0$, $Y_{m,0}=0$ almost surely.  If $k$ is nonintegral, the same two formulas with
$s$ replaced by $k$ hold respectively under the sharp convergence conditions
$m+2k>1$ and $m+2k>2$.
\end{theorem}

\begin{proof}
The coefficients $a_n^{(m,k)}$ are nonnegative.  Gauss' summation formula gives, because
$m+2k>0$,
\begin{equation}
 F_m(1)=\sum_{n\geq0}a_n^{(m,k)}
 =\frac{\Gamma(m)\Gamma(m+2k)}{\Gamma(m+k)^2}>0
 \label{eq:cueM-Gauss-normalizer}
\end{equation}
\cite[Gauss summation formula]{oldeDaalhuisDLMF15}.  This proves normalization and
\eqref{eq:cueM-Gauss-PGF}.  For integer $s$, $(-s)_n=0$ exactly when $n>s$; for noninteger
$k$ no Pochhammer factor vanishes.

In the nonterminating case, gamma notation gives
\[
 p_n^{(m,k)}=\frac{\Gamma(m)}{\Gamma(-k)^2F_m(1)}
 \frac{\Gamma(n-k)^2}{\Gamma(n+m)\Gamma(n+1)}.
\]
The standard gamma-ratio asymptotic
$\Gamma(n+a)/\Gamma(n+b)\sim n^{a-b}$ proves
\eqref{eq:cueM-Gauss-tail}.  Its power is strictly larger than $2$ exactly when
$m+2k>1$, and strictly larger than $3$ exactly when $m+2k>2$, explaining the two
nonintegral moment conditions.

Where the differentiated series converges at one,
\begin{align*}
 F_m'(1)&=\frac{k^2}{m}{}_2F_1(1-k,1-k;m+1;1),\\
 F_m''(1)&=\frac{k^2(k-1)^2}{m(m+1)}
 {}_2F_1(2-k,2-k;m+2;1).
\end{align*}
Applying Gauss' formula to these two values and dividing by
\eqref{eq:cueM-Gauss-normalizer} gives
\[
 \frac{F_m'(1)}{F_m(1)}=\frac{k^2}{m+2k-1},\qquad
 \frac{F_m''(1)}{F_m(1)}=
 \frac{k^2(k-1)^2}{(m+2k-2)(m+2k-1)}.
\]
They are respectively $\mathbb EY_{m,k}$ and
$\mathbb E[Y_{m,k}(Y_{m,k}-1)]$.  For positive integral $s$, the series terminates and the
same identities follow from Chu--Vandermonde without convergence restrictions; $s=0$ is the
separate point mass already stated.  This repairs the earlier version's unrestricted display
of the mean, without making any claim about the theorem's status outside that source.
\end{proof}

\paragraph{The fixed-length multiplier set is a ratio-lattice image.}
For a finite real slice $\beta^{(1)},\ldots,\beta^{(r)}\in(-1,1)$, set
\begin{equation}
 q_i=\frac{1+\beta^{(i)}}{1-\beta^{(i)}}>0,\qquad
 N_i=\begin{pmatrix}q_i&1\\0&1\end{pmatrix}.
 \label{eq:cueM-drift-matrices}
\end{equation}
The top-left coordinate is a monoid morphism because
\begin{equation}
 N_{i_1}\cdots N_{i_m}=
 \begin{pmatrix}
 \prod_{j=1}^{m}q_{i_j}&
 1+q_{i_1}+q_{i_1}q_{i_2}+\cdots+q_{i_1}\cdots q_{i_{m-1}}\\
 0&1
 \end{pmatrix}.
 \label{eq:cueM-drift-product}
\end{equation}

\begin{theorem}[Exact fixed-length multiplier rank]
\label{thm:cueM-multiplier-rank}
Let
\begin{equation}
 \mathcal M_m=\left\{q_{i_1}\cdots q_{i_m}:1\leq i_j\leq r\right\},
 \qquad
 H=\left\langle\frac{q_1}{q_r},\ldots,\frac{q_{r-1}}{q_r}\right\rangle
 \leq\mathbb R_{>0}^{\times},
 \label{eq:cueM-multiplier-set}
\end{equation}
and let $d=\operatorname{rank}_{\mathbb Z}H$.  Then
\begin{equation}
 |\mathcal M_m|=\Theta(m^d).
 \label{eq:cueM-multiplier-growth}
\end{equation}
More precisely,
\begin{equation}
 \mathcal M_m=q_r^m\left\{
 \prod_{i=1}^{r-1}(q_i/q_r)^{n_i}:
 n_i\in\mathbb Z_{\geq0},\ \sum_{i=1}^{r-1}n_i\leq m
 \right\}.
 \label{eq:cueM-multiplier-simplex-map}
\end{equation}
If $q_1,\ldots,q_r$ are multiplicatively independent, the displayed simplex map is
injective and
\begin{equation}
 |\mathcal M_m|=\binom{m+r-1}{r-1}.
 \label{eq:cueM-multiplier-independent}
\end{equation}
For arbitrary positive $q_i$,
\begin{equation}
 \operatorname{rank}_{\mathbb Z}\langle q_1,\ldots,q_r\rangle
 =\dim_{\mathbb Q}\operatorname{span}_{\mathbb Q}
 \{\log q_1,\ldots,\log q_r\};
 \label{eq:cueM-log-rank}
\end{equation}
no algebraicity hypothesis is needed.
\end{theorem}

\begin{proof}
A word is determined at the scalar level by its Parikh vector
$(n_1,\ldots,n_r)$ with $\sum_i n_i=m$.  Replacing
$n_r$ by $m-\sum_{i<r}n_i$ gives
\eqref{eq:cueM-multiplier-simplex-map}.  Since a finitely generated subgroup of
$\mathbb R_{>0}^{\times}$ is torsion-free, $H\cong\mathbb Z^d$.  Choose a basis of $H$.
The exponent vector of each generator $q_i/q_r$ is fixed, so every coordinate of every image
of the integer simplex is bounded in absolute value by $Cm$.  This gives $O(m^d)$ images.

Because the listed ratios generate a rank-$d$ group, some $d$ of their exponent vectors are
linearly independent over $\mathbb Q$, hence over $\mathbb Z$.  Restricting
\eqref{eq:cueM-multiplier-simplex-map} to nonnegative powers of those $d$ generators gives
an injective set of $\binom{m+d}{d}$ products.  This proves the lower bound, including the
case $d=0$.  If all $q_i$ are multiplicatively independent, equality of two fixed-length
products forces equality of every Parikh coordinate, giving
\eqref{eq:cueM-multiplier-independent}.

Finally, logarithm is an injective group homomorphism
$\mathbb R_{>0}^{\times}\to(\mathbb R,+)$ with exponential as inverse on its image.
Integral multiplicative relations among the $q_i$ are therefore exactly integral additive
relations among $\log q_i$; clearing denominators gives the same statement for rational
relations.  This proves \eqref{eq:cueM-log-rank}.  The earlier version counted with the rank
of $\langle q_i\rangle$; fixed word length removes the common $q_r^m$ direction, so the
ratio lattice $H$ is the exact object.
\end{proof}

\paragraph{The cubic trace polynomial.}
Apply the invertible affine coordinate change already proved in
\eqref{eq:n3-u-v-u-X-isomorphism}:
\begin{equation}
 X=9u-16v+1,\qquad v=\frac{9u+1-X}{16}.
 \label{eq:cueM-u-v-X-map}
\end{equation}
It has singleton fibres and deletes no coordinate.

\begin{theorem}[Generic splitting group of the cubic descended kernel]
\label{thm:cueM-Q3-S3}
Substitution of \eqref{eq:cueM-u-v-X-map} into \eqref{eq:cueM-Psi3} gives
\begin{equation}
 \Psi_3^\sharp\!\left(u,\frac{9u+1-X}{16}\right)
 =\frac5{16384}Q_3(u,X),
 \label{eq:cueM-Q3-transform}
\end{equation}
where
\begin{align}
 Q_3(u,X)={}&-7X^3+(549u+669)X^2
 -(6645u^2+22746u+8997)X\notag\\
 &+13783u^3+125949u^2+138933u+19087.
 \label{eq:cueM-Q3}
\end{align}
As a polynomial in $X$ over $\mathbb Q(u)$, $Q_3$ is irreducible and its splitting field has
Galois group $S_3$.
\end{theorem}

\begin{proof}
Equation \eqref{eq:cueM-Q3-transform} is direct coefficient collection after the displayed
affine substitution.  At $u=0$,
\[
 Q_3(0,X)=-7X^3+669X^2-8997X+19087.
\]
Modulo $11$, multiplication by the inverse of its nonzero leading scalar gives
\[
 f(X)=X^3+5X^2+3X-5.
\]
Its values at $0,1,\ldots,10$ are
\[
 6,4,7,10,8,7,2,10,4,1,7,
\]
so it has no root in $\mathbb F_{11}$ and is irreducible there.  Thus
$Q_3(0,X)$ is irreducible over $\mathbb Q$.  If $Q_3(u,X)$ factored over
$\mathbb Q(u)$, Gauss' lemma would give a factorization in $\mathbb Q[u][X]$ after units
were cleared.  Since the leading $X$ coefficient is the constant $-7$, the leading
coefficients of both positive-$X$-degree factors would be nonzero constants; specialization
at $u=0$ would retain both degrees and contradict the preceding irreducibility.  Therefore
$Q_3$ is irreducible over $\mathbb Q(u)$.

For a cubic $aX^3+bX^2+cX+d$, substitution in
$b^2c^2-4ac^3-4b^3d-27a^2d^2+18abcd$ gives
\begin{align}
 \operatorname{disc}_XQ_3
 =2^{18}3^3(&290275u^6+731094u^5+2234709u^4+2884436u^3\notag\\
             &+3453789u^2+1714038u+984779).
 \label{eq:cueM-Q3-discriminant}
\end{align}
This discriminant is not a square in $\mathbb Q(u)$.  Indeed, its nonzero specialization at
$u=0$ is $2^{18}3^3\cdot984779$; the final integer has digit sum $44$ and is not divisible
by $3$, so the $3$-adic valuation is the odd number $3$.  A rational-function square that is
regular and nonzero at $u=0$ specializes to a rational square there, which is impossible.

Irreducibility makes the Galois group a transitive subgroup of $S_3$.  The discriminant is a
square exactly when that group lies in $A_3$.  Here it does not, so the transitive group is
the full $S_3$.  The exact substitution, modular values, fraction-field irreducibility, and
discriminant expansion are independently replayed by the finite certificate below.
\end{proof}

\paragraph{Certificate and exact scope.}
The executable
\path{certificates/cue_mellin_programme/verify_cue_mellin_finite_core.py}
is recorded as \texttt{CERT-CUE-MELLIN-20260828-0001}.  It uses exact integer, rational,
Laurent-constant-term, and symbolic polynomial arithmetic.  Two independent installed
interpreters replay all 39 named checks: Python~3.13.9 with SymPy~1.13.1 and Python~3.13.1
with SymPy~1.13.3.  It verifies the $U(2)$ Weyl/terminating-${}_3F_2$ equality through order
seven; the exact $N=3$ recursive-state arrow and its squared-modulus identity; the trace
formula, exact degree, and leading coefficient through order eight; all four printed
$\Psi_s^\sharp$; the Gram obstruction; the affine cubic transform; fraction-field
irreducibility and the complete discriminant; and 42 terminating Gauss probability laws and
their first two factorial moments.

The analytic and structural statements are not delegated to that bounded computation:
\cref{thm:cueM-fixedN-support,thm:cueM-Hausdorff-Bernstein,%
thm:cueM-no-one-scale,thm:cueM-measure-peel,thm:cueM-transform-algebra,%
cor:cueM-U2-3F2,thm:cueM-U3-compact-measures,thm:cueM-U3-trace-descent,%
thm:cueM-Gauss-law,thm:cueM-multiplier-rank,thm:cueM-Q3-S3}
contain their standard proofs and exact dependencies.  The programme therefore supplies a
fixed-rank compact Mellin measure, a constructive integer-to-real-moment inverse, an exact
measure-valued derivative peel, a Laplace--Hurwitz--digamma transform carrier, complete
$U(2)$ moments, and a phase-retaining $U(3)$ trace descent.  The maps to the published CUE
derivative framework are \eqref{eq:cueM-GMS-specialization},
\eqref{eq:cueM-U3-first-state}, and \eqref{eq:cueM-U3-second-state}; the map to the separate
deck curve and its exact obstruction are already
\cref{thm:n3-deck-interface}.  None of these finite-rank or transform statements supplies a
map from an individual zeta zero to the CUE variables.  Any such further implication remains
a separately typed morphism and proof obligation.
